% English draft based on sections/05_emergent_stopping_layer_ja.tex
\section{The Emergent Stopping Layer $L_E$}

This chapter is the second center of the paper. \S{}3 showed that $L_F$ does not close by itself. This chapter shows how intelligence nevertheless stops: through the log layer $L_R$ and the emergent stopping layer $L_E$.

\begin{figure}[H]
\centering
\includegraphics[width=\linewidth]{fig3_mutual_constraint_mu_theta.pdf}
\caption{Mutual constraint $H_{ij}$ and threshold mobility $\mu_\theta$. Stopping regions are not generated by an isolated inference layer; they are projected through mutual constraint.}
\label{fig:part2-mutual-constraint}
\figurenote{$H_{ij}$ changes where stopping can occur. $\mu_\theta$ reorganizes the possible installation region of the stopping operator $E$.}{Stopping is not found inside $L_F$. It is projected and emerges as $L_E$ from the relation between $L_R$ and $H_{ij}$. Movable stopping is a condition of sustained freedom.}
\end{figure}

\begin{figure}[H]
\centering
\includegraphics[width=\linewidth]{fig4_three_layer_architecture.pdf}
\caption{The three-layer architecture of intelligence. Intelligence consists of $L_F$, $L_R$, and $L_E$; $\mu_\theta$ is the opening window through which stopping regions can be reconfigured.}
\label{fig:part2-three-layer}
\figurenote{$L_F$ cannot stop by itself. Its motion leaves traces in $L_R$, and recursive log motion under $H_{ij}$ gives rise to $L_E$.}{The figure joins the two centers of the paper: non-closure in \S{}3 and emergent stopping in \S{}5. The limit of $L_F$ becomes the condition for higher-layer formation.}
\end{figure}

\subsection{Problem Setting}

Real intelligence stops. Humans end a chain of reasoning, make judgments, and act. Since \S{}3 showed that stopping does not come from inside $L_F$, it must come from elsewhere. The question is what ``elsewhere'' means.

The answer is that stopping comes from a layer that is internal to the individual but externally dependent: the emergent stopping layer $L_E$.

\subsection{The Log Layer $L_R$}

The log layer $L_R$ is the layer of recursive reference to traces left by inference. These traces include passed states, attempted branches, failed paths, and points of divergence. $L_R$ is inside the individual in many systems, but outside the current operation of $F$.

The location of $L_R$ differs by system. In a slime mold, logs can be externalized into the environment as slime traces. In a Transformer, training data are frozen logs and context tokens are volatile logs. In sapiens-like intelligence, logs exist both individually and socially.

\subsection{Definition of $L_E$}

\textbf{Definition 5.1 (Emergent stopping layer).}
The emergent stopping layer $L_E$ is a stopping layer that dynamically emerges as a function of the recursive motion of $L_R$ and external mutual constraint $H_{ij}$:
\[
L_E=\Phi(L_R\text{ motion},H_{ij}).
\]
It is internal to the agent but externally dependent.

\begin{figure}[H]
\centering
\includegraphics[width=\linewidth]{fig5_three_examples_of_emergence.pdf}
\caption{Three examples of VED-defined emergence. Emergence is the formation of an upper attractor from lower non-closed motion under external mutual constraint $H_{ij}$.}
\label{fig:part2-emergence-examples}
\figurenote{Particle attractors, the life attractor $\Omega_{\text{life}}$, and the stopping layer $L_E$ can all be read as substitutions into $U=\Phi(L,H_{ij})$.}{Emergence is not mere complexity. It is the stabilization of non-closed motion as an attractor under mutual constraint, linking $L_E$ to the broader VED theory tree.}
\end{figure}

\subsection{Three Profiles}

Three systems clarify the variability of $L_R$, $\mu_\theta$, and $L_E$.

\textbf{Slime mold.} Its logs are environmental; its stopping is generated by direct environmental interaction; its $\mu_\theta$ is largely fixed.

\textbf{Transformer.} Its logs are internalized as frozen weights and volatile context; its stopping is externally injected through EOS tokens, task specifications, or decoding controls; it does not generate its own $L_E$ in the sense used here.

\textbf{Sapiens.} Logs are both individual and social; $L_E$ emerges internally and is socially synchronized; $\mu_\theta$ can be moved endogenously.

\subsection{Emergence Mechanism}

The emergence of $L_E$ proceeds in four stages:
\begin{enumerate}
\item $\mu_\theta$ causes region reconfiguration (B);
\item the reconfiguration projects a possible installation region for the stopping operator $E$;
\item the recursive motion of $L_R$ converges into the projected region;
\item the resulting stable stopping attractor is $L_E$.
\end{enumerate}

Schematically:
\[
\mu_\theta \to (B)\text{ reconfiguration} \to E\text{-region projection} \to L_R\text{ convergence} \to L_E .
\]

This mechanism shows that stopping is not the negation of freedom. Without projected stopping regions, inference becomes unsustainable in branching explosion and infinite regress. $L_E$ emerges not to restrict freedom, but to make free inference sustainable within a finite system.

\subsection{Projection}

The projection used here is the intelligence-layer realization of the projection primitive fixed in the consciousness paper. It means that a state or relation acquires unequal influence within a field. In this paper, the projected object is a possible stopping locus.

\subsection{External Stopping Operator}

The external stopping operator $E$ is the action of $L_E$ on an inferential state:
\[
E=E[L_E,x_n],\qquad E\notin\{F^k\}.
\]
It is external because it is not a member of the sequence of inference operations. It acts through $L_E$.

\subsection{$\Phi_{\text{ext}}\equiv\Phi_{\text{net}}$}

The external reference field $\Phi_{\text{ext}}$ is isomorphic to the inter-individual network of emergent stopping layers $\Phi_{\text{net}}$:
\[
\Phi_{\text{ext}}\equiv\Phi_{\text{net}}.
\]
Here $\Phi_{\text{ext}}$ should be read as an intelligence-layer projection of the IFGT potential vocabulary, not as the full IFGT field itself. It names the way a shared stopping structure appears as an external reference field from within an individual inference window.

\begin{figure}[H]
\centering
\includegraphics[width=\linewidth]{fig6_phi_ext_phi_net.pdf}
\caption{$\Phi_{\text{ext}} \equiv \Phi_{\text{net}}$. What appears to an individual as an external reference field is the network formed by mutual synchronization among individual $L_E$ layers.}
\label{fig:part2-phi-ext-net}
\figurenote{$\Phi_{\text{ext}}$ is the external appearance; $\Phi_{\text{net}}$ is the internal network structure of mutually synchronized $L_E$ layers.}{No transcendent field needs to be assumed. Norms, gods, truth, and monetary value can be read as appearances of distributed $L_E$ synchronization.}
\end{figure}

\subsection{Summary}

Intelligence persists through the simultaneous motion of $L_F$, $L_R$, and $L_E$ under externally dependent stopping. Sapiens-like intelligence adds endogenous $\mu_\theta$, allowing stopping loci to be reorganized rather than merely received.

\bigskip
